%% ma: L12-B16-0017
%% lop: 12
%% chuong: 5
%% bai: 16
%% dang: 12.16.10
%% loai: TN4
%% muc: TH
%% dap_an: "C"
%% nguon: "Ảnh giáo viên gửi 04/11/2026 (Huỳnh Văn Ánh, Chuyên đề V - Phương pháp tọa độ trong không gian), Đề test theo bài, Phần I Câu 8"
%% kien_thuc: "Bài 16 (góc giữa đường thẳng và mặt phẳng), Bài 7 (vectơ đơn vị của trục Ox), Bài 14 (vectơ pháp tuyến)"
%% trang_thai: chinh-thuc
%% ghi_chu: "Đáp án gốc C khớp (sympy): sin alpha = 0 và Ox song song (P) vì O không thuộc (P). Đã duyệt 04/11/2026."
\begin{ex}
Trong không gian với hệ tọa độ $Oxyz$, góc giữa đường thẳng $d$ chứa trục $Ox$ và mặt phẳng $(P)\colon 3y-z+5=0$ bằng
\choice{$90^\circ$}{$60^\circ$}{\True $0^\circ$}{$30^\circ$}
\loigiai{Đường thẳng $d$ chứa trục $Ox$ có vectơ chỉ phương $\vec u=(1;0;0)$; mặt phẳng $(P)$ có vectơ pháp tuyến $\vec n=(0;3;-1)$.

Gọi $\alpha$ là góc giữa $d$ và $(P)$: $\sin\alpha=\dfrac{|\vec u\cdot\vec n|}{|\vec u|\cdot|\vec n|}=0$, suy ra $\alpha=0^\circ$ (thật vậy $O\notin(P)$ nên $d\parallel(P)$). Chọn C.}
\end{ex}
